Fix an integer $n\ge 1$. Define a sequence $(a_k)_{k\ge 0}$ by
\[a_0:=n,\qquad a_1:=1,\]
and for each integer $k\ge 2$, define $a_k$ to be the least integer $>a_{k-1}$ such that
\[\gcd(n-a_k,\,n-a_i)=1\quad\text{for all }1\le i<k.
\]
Because $a_k$ may exceed $n$, the differences $n-a_k$ may be negative; throughout I interpret $\gcd(x,y):=\gcd(|x|,|y|)$.

Define (when the sequence is infinite)
\[S(n):=\sum_{k\ge 1} \frac1{a_k}.
\]
Main question: does $S(n)\to\infty$ as $n\to\infty$?

The problem also asks about restricted subseries, e.g. summing only those indices $k$ for which $n-a_k$ is divisible by some prime $\le a_k$, or the complementary set of indices.

\textbf{Edge case:} for $n=2$ the construction produces $a_1=1,a_2=2,a_3=3$ and then stops (no further $a_k$ exist), because the condition forces $|n-a_k|=1$ once $n-a_2=0$ appears. For $n\ge 3$, $a_k\neq n$ for all $k\ge 1$ (since $\gcd(n-n,n-1)=n-1>1$), so this obstruction does not occur.
